\documentclass[11pt,a4paper]{article}
\usepackage[T1]{fontenc}
\usepackage[utf8]{inputenc}
\usepackage{lmodern}
\usepackage{amsmath,amssymb,amsthm,mathtools,mathrsfs,bm}
\usepackage[margin=1in]{geometry}
\usepackage{microtype,booktabs,array}
\usepackage[hidelinks]{hyperref}
\usepackage[nameinlink,noabbrev]{cleveref}
\numberwithin{equation}{section}
\newtheorem{theorem}{Theorem}[section]
\newtheorem{lemma}[theorem]{Lemma}
\newtheorem{proposition}[theorem]{Proposition}
\newtheorem{corollary}[theorem]{Corollary}
\theoremstyle{definition}
\newtheorem{definition}[theorem]{Definition}
\newtheorem{remark}[theorem]{Remark}
\allowdisplaybreaks[1]
\title{An Infinite-Dimensional Smooth Kernel for an Exact Coupled Exponential Moment Operator}
\author{Author and affiliation to be confirmed}
\date{}
\begin{document}
\maketitle
\begin{abstract}
We study a coupled exponential moment operator whose rows combine a rapidly decaying exponential with a slowly accumulating exponential. A positive Bessel-kernel operator transforms these moments into a fast Laplace transform plus a dependent tail. For positive rates \(b_n\asymp n^2\) satisfying a separate uniform positive square-root gap, we construct on every compactly contained interior time interval an infinite-dimensional real space of smooth compactly supported corrections satisfying all coupled moments and an endpoint condition. The construction combines a positive-exponential biorthogonal estimate, M\"untz incompleteness, and a Fredholm finite-codimension argument. A conditional interpretation is given for a frozen linear odd heat--memory system; terminal temperature vanishes while all but finitely many terminal memory modes are nonzero for each nonzero constructed input. No nonlinear preparation or full-system inverse result follows.
\end{abstract}
\noindent\textbf{Keywords:} exponential moments; biorthogonal families; M\"untz theorem; compact operators; heat equations with memory.\\
\textbf{MSC 2020:} 44A10, 30E05, 93B05, 35R09.
\section{Introduction}\label{sec:intro}
The moment method reduces certain control problems for evolution equations to interpolation against exponentials. In one-dimensional parabolic problems, eigenvalues typically grow quadratically, and a biorthogonal family can be used to solve a suitable sequence of moment equations. The classical work of Fattorini and Russell \cite{FattoriniRussell1971} is a starting point for this method; for the precise positive-exponential estimate needed here we use Cannarsa, Martinez and Vancostenoble \cite{Cannarsa2020}. The latter provide quantitative bounds under square-root gap hypotheses. Their estimate concerns one separated family, not the coupled rows considered below.

Coupled evolution systems may exhibit a fast branch of eigenvalues and a second branch accumulating at a finite point. In the model below, after a time reversal the temperature-null condition pairs a fast Laplace value at \(-b_n\) with a slow value at \(d/b_n\). These values come from the same input. They are not independently selectable interpolation data, and the full family is not a single uniformly separated exponential family.

M\"untz theory \cite{Almira2007} implies incompleteness of the fast exponentials because \(\sum_n b_n^{-1}<\infty\). This observation alone does not prove incompleteness of the paired rows: a perturbation, even a compact one, can change a kernel. Similarly, dense range of a moment map concerns the adjoint kernel and cannot be substituted for injectivity of the map itself. Our argument instead represents the slow contribution as a dependent tail beginning at the terminal time. The gap between that time and an interior support interval yields an exponentially weighted compact map. A smooth right inverse for the fast moments and a Fredholm range intersection then provide exact simultaneous solutions without a smallness or invertibility assumption for the perturbation.

Heat equations with memory provide related, but not interchangeable, controllability questions. Guerrero and Imanuvilov \cite{Guerrero2013} use an observability argument to show failure of boundary null controllability for certain nonzero initial states in a different heat-memory equation. Halanay and Pandolfi \cite{HalanayPandolfi2015} study a Coleman--Gurtin equation, deriving approximate controllability and, under their stated hypotheses, failure of controllability to zero for some initial conditions. In particular, their treatment highlights the need to specify when terminal evaluation of a boundary-controlled solution is meaningful. Ivanov and Pandolfi \cite{IvanovPandolfi2009} study controllability to rest, a target stronger than vanishing temperature at one time; we have not inspected their original proof. None of these comparisons by itself identifies or excludes an earlier version of the present exact coupled-kernel result.

We formulate a frozen linear odd model to motivate the rows, then prove an abstract theorem independently of that interpretation. The theorem constructs an infinite-dimensional space of interior-supported smooth corrections, imposes an endpoint condition on the transformed input, and solves every coupled equation exactly. The PDE consequence concerns smooth compatible boundary inputs in the stated linear model, not preparation of states in a nonlinear experiment. The literature comparison is targeted rather than exhaustive; independent proof and priority review remain necessary.

\begin{table}[htbp]
\centering\small
\begin{tabular}{@{}p{.20\linewidth}p{.36\linewidth}p{.36\linewidth}@{}}
\toprule
Source & Problem and method & Relation to this paper \\
\midrule
\cite{Cannarsa2020} & Separated exponentials; quantitative biorthogonal estimates. & Supplies the fast right inverse, not coupled-row incompleteness. \\
\cite{Almira2007} & Uniform approximation by powers; M\"untz criteria. & Supplies fast-only null seeds, not the dependent-tail conclusion. \\
\cite{Guerrero2013} & Boundary null controllability from nonzero initial states; observability. & Different target and initial data from zero-start paired nulling. \\
\cite{HalanayPandolfi2015} & Coleman--Gurtin reachability and non-null-controllability via heat-equation properties. & Different memory operator and reachability question. \\
\cite{IvanovPandolfi2009} & Controllability to rest in a heat-memory equation. & Stronger target; original proof not inspected. \\
\bottomrule
\end{tabular}
\caption{Targeted comparison; differences do not establish novelty or priority.}\label{tab:literature}
\end{table}

\section{Governing model and preliminaries}\label{sec:model}
For the interpretation, consider real fields on \(0<x<L\), \(-T<s<0\), satisfying the frozen equations
\begin{equation}\label{eq:pde}
M Y_s=\kappa Y_{xx}+B_0Y+SP,\qquad P_s=-\mu P+\gamma Y,
\end{equation}
with \(M,\kappa,L,\mu>0\), \(d=\gamma S/M>0\), and real constant \(B_0,S,\gamma\). Initially \(Y(\cdot,-T)=P(\cdot,-T)=0\); the prescribed boundary values are \(Y(0,s)=a(s)\), \(Y(L,s)=-a(s)\). We work in the antisymmetric subspace about \(x=L/2\). This is a linear frozen model, not an assertion about nonlinear reachability or prescribed physical parameter values.

For the interpretation below take a smooth control with \(a(-T)=a(0)=0\), and use a smooth compatible boundary lift \(Y=Z+a(s)(1-2x/L)\). The homogeneous-Dirichlet heat generator on \(L^2(0,L)\), coupled to the bounded memory ODE, has a bounded off-diagonal perturbation. The lift supplies continuous \(L^2\)-valued forcing, so standard semigroup variation of constants gives a continuous \(L^2\)-valued pair. The basis for the odd subspace is
\begin{equation}\label{eq:basis}
\phi_n(x)=\sqrt{2/L}\sin(q_nx),\qquad q_n=2\pi n/L,\quad n\ge1.
\end{equation}
All terminal identities below are understood in this solution class, not as a general endpoint trace theorem for arbitrary \(L^2\) boundary data. The physical relation \(\mu B_0+\gamma S=0\), when imposed, is an additional condition; it is not assumed for the abstract theorem.

\section{Modal decomposition and exact moments}\label{sec:modal}
Set \(Y_n=\langle Y,\phi_n\rangle\), \(P_n=\langle P,\phi_n\rangle\). Integration by parts and \(\phi_n'(0)=\phi_n'(L)=\sqrt{2/L}q_n\) yield
\begin{equation}\label{eq:ibp}
\int_0^L Y_{xx}\phi_n\,dx=-q_n^2Y_n+2a(s)\sqrt{2/L}\,q_n.
\end{equation}
Consequently, with \(A_n=(\kappa q_n^2-B_0)/M\) and \(C_n=2\kappa\sqrt{2/L}\,q_n/M\),
\begin{equation}\label{eq:modalode}
\frac d{ds}\binom{Y_n}{P_n}=D_n\binom{Y_n}{P_n}+\binom{C_na}{0},\qquad
D_n=\begin{pmatrix}-A_n&S/M\\\gamma&-\mu\end{pmatrix}.
\end{equation}
The characteristic equation is \((\lambda+A_n)(\lambda+\mu)-d=0\). Define
\begin{equation}\label{eq:roots}
b_n=\frac{A_n-\mu+\sqrt{(A_n-\mu)^2+4d}}2>0,\qquad \delta_n=\frac d{b_n}>0.
\end{equation}
Then the roots are \(-\mu-b_n\), \(-\mu+\delta_n\), and \(b_n-\delta_n=A_n-\mu\). The root \(b_n\) increases strictly with \(A_n\). Writing \(\alpha=\kappa(2\pi/L)^2/M>0\), \(\beta=-B_0/M-\mu\), gives \(b_n=\alpha n^2+\beta+O(n^{-2})\); its square-root gaps tend to \(\sqrt\alpha\). Strict increase and this limit give a positive global gap. If \(\mu B_0+\gamma S=0\), the determinant \(\mu A_n-d=\mu\kappa q_n^2/M\) and the trace are positive and negative respectively, so both roots are negative.

For \(f(t)=a(-t)\) and \(F(z)=\int_0^T e^{zt}f(t)dt\), variation of constants gives \(\binom{Y_n(0)}{P_n(0)}=C_n\int_0^T e^{tD_n}\binom10 f(t)dt\). The two spectral projectors give
\begin{equation}\label{eq:projector}
e^{tD_n}\binom10=\frac1{b_n+\delta_n}
\binom{\delta_ne^{(-\mu+\delta_n)t}+b_ne^{(-\mu-b_n)t}}
{\gamma(e^{(-\mu+\delta_n)t}-e^{(-\mu-b_n)t})}.
\end{equation}
Thus
\begin{align}
Y_n(0)&=\frac{C_n}{b_n+\delta_n}\left[\delta_nF(-\mu+\delta_n)+b_nF(-\mu-b_n)\right],\label{eq:Ymodal}\\
P_n(0)&=\frac{C_n\gamma}{b_n+\delta_n}\left[F(-\mu+\delta_n)-F(-\mu-b_n)\right].\label{eq:Pmodal}
\end{align}
If \(f(t)=e^{\mu t}t v(t)\), the zero-temperature condition is precisely
\begin{equation}\label{eq:paired}
\int_0^T t v(t)\left(e^{-b_nt}+\frac d{b_n^2}e^{\delta_nt}\right)dt=0,\qquad n\ge1.
\end{equation}
The coefficient of the slow term has a plus sign and is \(d/b_n^2\). Its values are tied to the same \(v\) as the fast values.

\section{Bessel-kernel operator and functional setting}\label{sec:setting}
From here the theorem is abstract: fix \(T,d,\mu>0\) and an increasing positive sequence satisfying
\begin{equation}\label{eq:growth}
b_n\asymp n^2,\qquad \inf_n(\sqrt{b_{n+1}}-\sqrt{b_n})>0.
\end{equation}
The gap is a separate hypothesis: growth order alone permits clustered pairs. All spaces are real. Put \(H=L^2((0,T),t\,dt)\) and
\begin{equation}\label{eq:kernel}
H_d(t,s)=\sum_{j=0}^\infty\frac{(dts)^j}{j!(j+1)!},\qquad
(\mathsf Tv)(t)=d\int_0^T H_d(t,s)v(s)s\,ds.
\end{equation}
\begin{lemma}\label{lem:positive}
The operator \(\mathsf T\) is compact, self-adjoint and nonnegative on \(H\). Its quadratic form is
\begin{equation}\label{eq:positivity}
\langle\mathsf Tv,v\rangle_H=d\sum_{j=0}^\infty\frac{d^j}{j!(j+1)!}\left(\int_0^T s^{j+1}v(s)ds\right)^2.
\end{equation}
The inverse \(Q=(I+\mathsf T)^{-1}\) exists with \(\|Q\|\le1\). For \(v\in H\), \(\mathsf Tv\) extends to an entire function of \(t\); consequently \(Qk\) is smooth on \([0,T]\) for smooth \(k\).
\end{lemma}
\begin{proof}
The series converges uniformly on compact sets in \((t,s)\), since \(j!(j+1)!\ge(j!)^2\). Thus the symmetric kernel is bounded on \([0,T]^2\) and belongs to \(L^2((0,T)^2,ts\,dt\,ds)\), proving Hilbert--Schmidt compactness and self-adjointness. Uniform convergence and Cauchy--Schwarz justify integration termwise, giving \eqref{eq:positivity}. Since \(I+\mathsf T\) is bounded below by one and self-adjoint, its range is closed and has zero orthogonal complement; it is onto and \(\|Q\|\le1\). Uniform convergence of the integral series on complex compact sets proves the entire extension. Finally \(Qk=k-\mathsf TQk\) almost everywhere, giving the asserted smooth representative.
\end{proof}
For \(k\in H\) define the correction moments
\begin{equation}\label{eq:B}
(\mathsf Bk)_n=\int_0^T te^{-b_nt}k(t)dt+
\int_T^\infty te^{-b_nt}(\mathsf TQk)(t)dt.
\end{equation}
The second integral is finite by the growth bound proved in \cref{lem:tail}. The identity
\begin{equation}\label{eq:laplacekernel}
\int_0^\infty te^{-bt}H_d(t,s)dt=b^{-2}e^{ds/b},\qquad b>0,\quad 0\le s\le T,
\end{equation}
follows by termwise integration: the integral of \(t^{j+1}e^{-bt}\) is \((j+1)!b^{-j-2}\). Nonnegative terms and then Cauchy--Schwarz justify the interchange. If \(v=Qk\), then \(k=v+\mathsf Tv\), and \eqref{eq:laplacekernel} gives
\begin{equation}\label{eq:exactidentity}
(\mathsf Bk)_n=\int_0^T t v(t)\left(e^{-b_nt}+d b_n^{-2}e^{dt/b_n}\right)dt.
\end{equation}
Thus the tail is not an independent unknown: it is determined by \(k\) through \(Q\).

\section{Main theorem}\label{sec:theorem}
For smooth \(k\), use the representative of \(Qk\) in \cref{lem:positive}, and set
\begin{equation}\label{eq:adm}
\mathcal K_{\rm adm}=\{k\in C^\infty[0,T]: k^{(j)}(0)=0\ (j\ge0),\ (Qk)(T)=0\}.
\end{equation}
\begin{theorem}\label{thm:main}
Assume \eqref{eq:growth}. For every open interval \(I\) with \(\overline I\subset(0,T)\), the real vector space
\begin{equation}\label{eq:kernelspace}
\{k\in C_c^\infty(I):\mathsf Bk=0,\ (Qk)(T)=0\}
\end{equation}
is infinite-dimensional. In particular, \(\ker(\mathsf B|_{\mathcal K_{\rm adm}})\) is infinite-dimensional.
\end{theorem}
The space includes zero; the assertion provides infinitely many linearly independent nonzero corrections. Its proof is given in \cref{sec:biorth,sec:fredholm,sec:endpoint}. It does not assert a numerical formula or control cost.

\section{Biorthogonal construction and smooth right inverse}\label{sec:biorth}
Fix \(I\Subset(0,T)\). Choose \(a,\ell,\varepsilon,\eta>0\) such that \([a,a+\ell+\varepsilon]\subset I\) and \(a+\varepsilon<T-\eta\). Let \(X=L^2(I,t\,dt)\), embedded in \(H\) by zero extension, and define
\begin{equation}\label{eq:EY}
(Ex)_n=\int_I te^{-b_nt}x(t)dt,\qquad
Y=\left\{y:\sum_n e^{2(T-\eta)b_n}|y_n|^2<\infty\right\}.
\end{equation}
The space \(Y\) has its indicated weighted \(\ell^2\) norm. The map \(E:X\to\ell^2\) is bounded: the sum of squares of its coordinate-functional norms is bounded by \(\sum_n\int_I te^{-2b_nt}dt<\infty\).

\begin{proposition}\label{prop:rightinverse}
There is a bounded linear \(S:Y\to X\) whose values have a common compact support in \(I\), belong to \(C_c^\infty(I)\), and satisfy \(ESy=y\) coordinatewise. Its defining series converges in every \(C^j\) norm.
\end{proposition}
\begin{proof}
Let \(g_0>0\) be the infimum in \eqref{eq:growth}. Theorem 2.1 of \cite{Cannarsa2020} applies to positive exponentials \(e^{b_nu}\) on \((0,\ell)\): take its global gap \(g_0/2\), asymptotic gap \(g_0\), and \(N_*=1\). Then \(M_*=0\), and the remaining source factor is bounded in \(n\) because \(b_1>0\). After taking square roots and fixing the time and gaps, there are real \(\sigma_n\in L^2(0,\ell)\) with
\begin{equation}\label{eq:sigma}
\int_0^\ell e^{b_mu}\sigma_n(u)du=\delta_{mn},\qquad
\|\sigma_n\|_2\le C e^{-\ell b_n+C\sqrt{b_n}}.
\end{equation}
Here \(\delta_{mn}\) is the Kronecker symbol, not the spectral parameter \(\delta_n\). Set \(p_n(u)=e^{\ell b_n}\sigma_n(\ell-u)\) on \((0,\ell)\), extended by zero. Changing variables gives
\begin{equation}\label{eq:pn}
\int_{\mathbb R}e^{-b_mu}p_n(u)du=\delta_{mn},\qquad \|p_n\|_2\le C e^{C\sqrt{b_n}}.
\end{equation}
Choose a nonzero nonnegative \(\rho\in C_c^\infty(0,\varepsilon)\), and put \(\rho_n=\int e^{-b_nu}\rho(u)du\). Then \(\rho_n\ge e^{-\varepsilon b_n}\int\rho>0\). Define
\begin{equation}\label{eq:psi}
\psi_n(t)=\frac{e^{ab_n}}{t\rho_n}[p_n(\cdot-a)*\rho](t).
\end{equation}
The convolution has a common compact support inside \(I\), where division by \(t\) is smooth. Factoring its Laplace transform gives \(\int_I te^{-b_mt}\psi_n(t)dt=\delta_{mn}\): on the diagonal the factors \(e^{ab_n}\), \(e^{-ab_m}\), \(\rho_m\) and \(\rho_n^{-1}\) cancel; off the diagonal \eqref{eq:pn} vanishes. By Cauchy--Schwarz on each convolution derivative and Leibniz' rule, for every \(j\ge0\),
\begin{equation}\label{eq:psibound}
\|\psi_n\|_{C^j}\le C_j e^{(a+\varepsilon)b_n+C\sqrt{b_n}}.
\end{equation}
Define \(Sy=\sum_n y_n\psi_n\). With \(c=T-\eta-a-\varepsilon>0\), Cauchy--Schwarz bounds the sum of \(C^j\) norms by \(C_j\|y\|_Y(\sum_n e^{-2cb_n+2C\sqrt{b_n}})^{1/2}\), which is finite using \(b_n\asymp n^2\). Hence the derivative series is uniformly absolutely convergent for every \(j\), with common compact support. Termwise integration proves \(ESy=y\). The right inverse is only asserted on \(Y\), not on all of \(\ell^2\).
\end{proof}

\section{Compact tail and Fredholm intersection}\label{sec:fredholm}
\begin{lemma}\label{lem:tail}
Define \(R:X\to Y\) by
\begin{equation}\label{eq:R}
(Rx)_n=\int_T^\infty te^{-b_nt}(\mathsf TQx)(t)dt.
\end{equation}
This map is Hilbert--Schmidt. Hence \(SR:X\to X\) is compact.
\end{lemma}
\begin{proof}
For \(z\ge0\), \(\sum_{j\ge0}z^j/(j!)^2\le e^{2\sqrt z}\), since the sum is dominated by the square of \(\sum_j(\sqrt z)^j/j!\). Thus, for \(t\ge T\),
\begin{equation}\label{eq:tailgrowth}
|(\mathsf TQx)(t)|\le d\|Qx\|_H\left(\int_0^T s|H_d(t,s)|^2ds\right)^{1/2}
\le C e^{2\sqrt{dTt}}\|x\|_X.
\end{equation}
Writing \(t=T+u\), and using \(b_n\ge b_1>0\),
\[
\int_T^\infty te^{-b_nt}e^{2\sqrt{dTt}}dt
\le e^{-Tb_n}\int_0^\infty(T+u)e^{-b_1u+2\sqrt{dT(T+u)}}du
\le C e^{-Tb_n}.
\]
The final integral is finite because its exponent is \(-b_1u+O(\sqrt u)\). The weighted coordinate functionals of \(R\) therefore have squared norms bounded by \(C^2e^{-2\eta b_n}\), whose sum is finite. For a map from a Hilbert space into a weighted \(\ell^2\) space this is exactly the Hilbert--Schmidt criterion. Composition with bounded \(S\) is compact.
\end{proof}
\begin{proposition}\label{prop:fredholm}
There is an infinite-dimensional subspace of \(C_c^\infty(I)\) on which \(E+R=0\) coordinatewise.
\end{proposition}
\begin{proof}
Let \(A=I_X+SR\). The Fredholm alternative \cite{Howard2020} gives closed range of finite codimension; neither invertibility of \(A\) nor smallness of \(SR\) is needed. We first show that
\[
N=\{h\in C_c^\infty(I):Eh=0\}
\]
is infinite-dimensional. Choose a closed interval \(J\) in the interior of \(I\). By \(b_n\asymp n^2\), \(\sum_n b_n^{-1}<\infty\). The full M\"untz theorem away from zero, \cite[Theorem~29]{Almira2007}, applied after \(z=e^{-t}\), shows that the fast exponentials are not dense in \(C(J)\). Hahn--Banach and the Riesz representation theorem \cite{RieszNotes} provide a nonzero finite signed measure \(\nu\) on \(J\) annihilating all of them.

Choose a nonzero compactly supported smooth bump \(\rho\) so small that \(J+\operatorname{supp}\rho\Subset I\). The function \(w=\nu*\rho\) is smooth and compactly supported, and its Laplace moments factor into those of \(\nu\) and \(\rho\), hence vanish. It is nonzero: the Fourier transforms of both compactly supported factors are nonzero entire functions, and their product is not identically zero. Then \(h=w/t\) is smooth, compactly supported and belongs to \(N\). Repeating the construction in pairwise disjoint subintervals of \(I\) gives linearly independent \(h\)'s.

The quotient \(X/\operatorname{ran}A\) is finite-dimensional. Therefore the kernel of the quotient map restricted to \(N\), namely \(N\cap\operatorname{ran}A\), is infinite-dimensional. Choose linearly independent \(h_i\) there and solve \(Ax_i=h_i\) in \(X\). Since \(x_i=h_i-SRx_i\) as an \(X\)-identity, \(x_i\) has a \(C_c^\infty(I)\) representative. Moreover, \(Ex_i=Eh_i-ESRx_i=-Rx_i\), because \(Rx_i\in Y\) and \(ES=I_Y\). The \(x_i\) are independent: applying \(A\) sends any linear dependence to the corresponding dependence among the \(h_i\). Their span gives the proposition. No finite-mode approximation is used.
\end{proof}

\section{Endpoint constraint and proof of the theorem}\label{sec:endpoint}
\begin{lemma}\label{lem:endpoint}
The functional
\begin{equation}\label{eq:endpoint}
L(x)=-d\int_0^T H_d(T,s)(Qx)(s)s\,ds,\qquad x\in X,
\end{equation}
is bounded. If \(x\in C_c^\infty(I)\), then \(L(x)=(Qx)(T)\), with the smooth representative.
\end{lemma}
\begin{proof}
The kernel row \(dH_d(T,\cdot)\) belongs to \(H\) and \(Q\) is bounded, so Cauchy--Schwarz gives \(|L(x)|\le C\|x\|_X\). Point evaluation is not defined for an arbitrary element of \(H\). However, the zero extension of \(x\in X\) vanishes almost everywhere near \(T\), and \(Qx=x-\mathsf TQx\) has the canonical representative \(-\mathsf TQx\) there. For smooth compactly supported \(x\), this agrees with the smooth representative on all of \([0,T]\), giving the claimed endpoint value.
\end{proof}
\begin{proof}[Proof of \cref{thm:main}]
Let \(V\subset C_c^\infty(I)\) be the infinite-dimensional space from \cref{prop:fredholm}. For every \(k\in V\), \((E+R)k=\mathsf Bk=0\), by the zero extension in \eqref{eq:B}. The restriction of one real linear functional \(L\) to \(V\) has a kernel of codimension at most one. Hence \(V\cap\ker L\) is infinite-dimensional, and \cref{lem:endpoint} gives \((Qk)(T)=0\) for each of its members. Such \(k\) also vanishes near zero with all derivatives, so it lies in \(\mathcal K_{\rm adm}\). This proves the theorem.
\end{proof}

\section{Conditional terminal-state interpretation}\label{sec:interpretation}
Take \(b_n\) from \eqref{eq:roots}, and a nonzero \(k\) in \eqref{eq:kernelspace}. Put \(v=Qk\), \(f(t)=e^{\mu t}tv(t)\), \(a(s)=f(-s)\). Then \(v\) is smooth, \(f(0)=0\), and \(f(T)=0\) by the endpoint condition. Under the smooth-lift solution framework in \cref{sec:model}, \eqref{eq:exactidentity} and \eqref{eq:Ymodal} imply
\begin{equation}\label{eq:temperature}
Y_n(0)=\frac{C_nb_n}{b_n+\delta_n}(\mathsf Bk)_n=0.
\end{equation}
All coefficients in the odd basis vanish, so the odd terminal temperature is zero in \(L^2(0,L)\). The zero-temperature relation in \eqref{eq:Ymodal} gives \(F(-\mu-b_n)=-d b_n^{-2}F(-\mu+\delta_n)\). Substitution in \eqref{eq:Pmodal} yields
\begin{equation}\label{eq:memory}
P_n(0)=\frac{C_n\gamma}{b_n}F(-\mu+\delta_n)
=\frac{C_n\gamma}{b_n}\int_0^T te^{\delta_nt}v(t)dt.
\end{equation}
No terminal condition \(P(0)=0\) is imposed.

Because \(Q\) is injective, \(v\ne0\). Each weight in \eqref{eq:paired} is positive on \((0,T)\), so the real continuous functions \(v\) and \(f\) must change sign. Since \(k\) vanishes near both endpoints, \(v=-\mathsf Tv\) there, and \(f\) is analytic near each endpoint. If \(f\) vanished on \((0,\epsilon)\), then \(\mathsf Tv\) would vanish there and hence everywhere by analyticity. By \eqref{eq:positivity}, all weighted moments of \(v\) would vanish; polynomial density in \(H\) would imply \(v=0\), a contradiction. Thus, in physical time, the control has no terminal switch-off gap.

The entire function \(G(z)=\int_0^T te^{zt}v(t)dt\) is nonzero: otherwise all weighted polynomial moments vanish and \(v=0\). Since \(\delta_n\to0\), its zeros cannot contain infinitely many \(\delta_n\) unless \(G\equiv0\). From \eqref{eq:memory}, \(C_n\ne0\) and \(\gamma\ne0\), so all but finitely many terminal memory coefficients are nonzero. Finally, \(f\) cannot be real analytic on a connected neighborhood of \([0,T]\): the quotient \(e^{-\mu t}f(t)/t\) has an analytic removable singularity at zero, making \(k=v+\mathsf Tv\) analytic and identically zero because it vanishes near zero. These are assertions about this linear odd solution class only.

\section{Discussion and limitations}\label{sec:discussion}
The construction combines two tools that address different parts of the problem: M\"untz incompleteness supplies fast-null seeds, while quantitative biorthogonality solves fast interpolation for the tail's rapidly decaying coordinates. The Fredholm step is applied to \(I+SR\) on the interior-support space, not to the original coupled moment map. Its range is used exactly, not replaced by its closure. The construction is existential and gives no uniform quantitative control cost or prescribed Gevrey order.

The asymptotic condition and the uniform square-root gap are both stated because the former alone does not imply the latter. The precise external estimate and its specialization require independent checking; our use of it is not a new proof of the cited result. The PDE interpretation additionally requires the specified frozen coefficients, zero initial odd fields, compatible boundary values and the smooth-lift solution class. It is not a nonlinear state-preparation theorem, a controllability-to-rest theorem, a proof of physical reachability, or an inverse result for the full system. The literature review is not exhaustive and establishes no chronology or novelty. The proof and model interpretation have not been independently certified.

\section{Conclusion}
Under quadratic spectral growth and a separate positive square-root gap, the dependent-tail representation and a weighted smooth right inverse yield an infinite-dimensional kernel of the exact coupled moment map, even after the endpoint constraint. In the stated frozen linear odd model these inputs give zero terminal temperature and generally nonzero terminal memory. Independent mathematical, model and literature review remain necessary before any submission or priority claim.

\bibliographystyle{unsrt}
\bibliography{references}
\end{document}
